\documentclass[11pt,a4paper]{article}

\usepackage{amsmath}
\usepackage{amssymb}
\usepackage{amsthm}
\usepackage{graphicx}
\usepackage{hyperref}
\usepackage{xcolor}
\usepackage{enumitem}

% Define colors used in the original
\definecolor{cyan}{RGB}{0,255,255}
\definecolor{limegreen}{RGB}{50,205,50}
\definecolor{orange}{RGB}{255,128,0}

\title{Sequential Monte Carlo and Improved Auxiliary Particle Filters}
\author{Nicola Branchini}
\date{}

\begin{document}

\maketitle

\begin{abstract}
In this post, my aims are: introduce Bayesian inference in state space models; introduce approximate inference using importance sampling in state space models, presenting and comparing different derivations found in the literature; and describe the Auxiliary Particle Filter, its diverse interpretations, and the recent Improved Auxiliary Particle Filter by Elvira et al.
\end{abstract}

\tableofcontents

\section{Brief Introduction to Sequential Inference}

In Bayesian inference we want to update our beliefs on the state of some random variables, which could represent parameters of a parametric statistical model or represent some unobserved data generating process. The probabilistic evolution of a dynamical system is often called a \textit{state space model}.

Let the (unknown) state of the system at time $t$ be the vector valued random variable $\mathbf{s}_t$. We observe this state through a (noisy) measurement $\mathbf{v}_t$ (where v stands for visible).

The state $\mathbf{s}_t$ is sampled from some density conditional on the previous timestep:

\begin{equation}
\label{eq:transition}
\textcolor{cyan}{\text{Transition density}}: \qquad \mathbf{s}_t \sim f(\mathbf{s}_t \mid \mathbf{s}_{t-1})
\end{equation}

The observation $\mathbf{v}_t$ is sampled according to the current state:

\begin{equation}
\label{eq:observation}
\textcolor{limegreen}{\text{Observation density}}: \qquad \mathbf{v}_t \sim g(\mathbf{v}_t \mid \mathbf{s}_t)
\end{equation}

This collection of random variables and densities defines the state space model completely. When the transition and observation densities are linear with additive Gaussian noise, the model is called a \textit{Linear Dynamical System} (LDS). When variables are discrete, it is often called a \textit{Hidden Markov Model} (HMM).

There are several tasks we can perform on the state space model:

\begin{itemize}
\item \textbf{Filtering}: The target distributions are $p(\mathbf{s}_t \mid \mathbf{v}_{1:t})$, for $t = 1, \ldots, T$. This represents what we have learnt about the system's state at time $t$, after observations up to $t$.

\item \textbf{Smoothing}: The target distributions are $p(\mathbf{s}_t \mid \mathbf{v}_{1:T})$, for $t = 1, \ldots, T$. This represents what we have learnt about the system's state after observing the complete sequence of measurements.

\item \textbf{Parameter Estimation}: The target distribution is $p(\boldsymbol{\theta} \mid \mathbf{v}_{1:T}) = \int p(\mathbf{s}_{0:T}, \boldsymbol{\theta} \mid \mathbf{v}_{1:T}) \, d\mathbf{s}_{0:T}$, where $\boldsymbol{\theta}$ represents all parameters of any parametric densities in the state space model.
\end{itemize}

\subsection{General Bayesian Filtering}

\subsubsection{Notation}

\begin{itemize}[leftmargin=*]
\item The notation $\mathbf{v}_{1:t}$ means a collection of vectors $(\mathbf{v}_1, \mathbf{v}_2, \dots, \mathbf{v}_t)$
\item Therefore, $p(\mathbf{v}_{1:t})$ is a joint distribution: $p(\mathbf{v}_1, \mathbf{v}_2, \dots, \mathbf{v}_t)$
\item Integrating $\int p(\mathbf{x}_{1:t}) \, d\mathbf{x}_{i:j}$ means $\underbrace{\int \cdots \int}_{j-i+1} p(\mathbf{x}_{1:t}) \, d\mathbf{x}_i \, d\mathbf{x}_{i+1} \cdots d\mathbf{x}_j$
\item The symbol $:=$ denotes a definition
\end{itemize}

Let's derive the filtering distribution. Recall that the aim is to compute $p(\mathbf{s}_t \mid \mathbf{v}_{1:t})$. Apply Bayes rule:

\begin{equation}
\label{eq:bayes_filter}
p(\mathbf{s}_t \mid \mathbf{v}_{1:t}) = \frac{\overbrace{p(\mathbf{v}_t \mid \mathbf{s}_t)}^{\mathbf{v}_t \text{ only dep. on } \mathbf{s}_t} \, p(\mathbf{s}_t \mid \mathbf{v}_{1:t-1})}{p(\mathbf{v}_t \mid \mathbf{v}_{1:t-1})} = \frac{g(\mathbf{v}_t \mid \mathbf{s}_t) \, p(\mathbf{s}_t \mid \mathbf{v}_{1:t-1})}{p(\mathbf{v}_t \mid \mathbf{v}_{1:t-1})}
\end{equation}

We need to compute the term $p(\mathbf{s}_t \mid \mathbf{v}_{1:t-1})$, which is a marginal:

\begin{equation}
p(\mathbf{s}_t \mid \mathbf{v}_{1:t-1}) = \int p(\mathbf{s}_t, \mathbf{s}_{t-1} \mid \mathbf{v}_{1:t-1}) \, d\mathbf{s}_{t-1}
\end{equation}

Continuing, we split the joint with the product rule:

\begin{align}
p(\mathbf{s}_t \mid \mathbf{v}_{1:t-1}) &= \int p(\mathbf{s}_t \mid \mathbf{s}_{t-1}) \, p(\mathbf{s}_{t-1} \mid \mathbf{v}_{1:t-1}) \, d\mathbf{s}_{t-1} \nonumber \\
&= \int f(\mathbf{s}_t \mid \mathbf{s}_{t-1}) \, p(\mathbf{s}_{t-1} \mid \mathbf{v}_{1:t-1}) \, d\mathbf{s}_{t-1}
\label{eq:prediction}
\end{align}

The term $p(\mathbf{s}_{t-1} \mid \mathbf{v}_{1:t-1})$ is the filtering distribution at $t-1$, which we have already computed recursively.

The step requiring computation of $p(\mathbf{s}_t \mid \mathbf{v}_{1:t-1})$ is called \textit{prediction}, and the step $p(\mathbf{s}_t \mid \mathbf{v}_{1:t}) \propto g(\mathbf{v}_t \mid \mathbf{s}_t) \cdot p(\mathbf{s}_t \mid \mathbf{v}_{1:t-1})$ is called \textit{correction}.

In an LDS all computations have closed form solutions, yielding the \textit{Kalman Filter}. For discrete valued random variables with small dimensionalities, we can use the \textit{Forward-Backward} algorithm for HMMs. When variables are non-Gaussian and/or densities are nonlinear, we need approximate inference.

\subsection{Recursive Formulations}

There is another way to derive generic computation steps. Consider the sequential estimation of:

\begin{equation}
\label{eq:joint_posterior}
p(\mathbf{s}_{1:t} \mid \mathbf{v}_{1:t}) \propto p(\mathbf{s}_{1:t}, \mathbf{v}_{1:t})
\end{equation}

Applying Bayes' rule gives the recursive relationship:

\begin{equation}
\label{eq:recursive_joint}
p(\mathbf{s}_{1:t}, \mathbf{v}_{1:t}) = p(\mathbf{s}_{1:t-1}, \mathbf{v}_{1:t-1}) \, f(\mathbf{s}_t \mid \mathbf{s}_{t-1}) \, g(\mathbf{v}_t \mid \mathbf{s}_t)
\end{equation}

For the sequential estimation of $p(\mathbf{s}_{1:t} \mid \mathbf{v}_{1:t})$:

\begin{align}
p(\mathbf{s}_{1:t} \mid \mathbf{v}_{1:t}) &= \frac{p(\mathbf{s}_{1:t}, \mathbf{v}_{1:t})}{p(\mathbf{v}_{1:t})} \nonumber \\
&= \frac{p(\mathbf{s}_{1:t-1}, \mathbf{v}_{1:t-1}) \, f(\mathbf{s}_t \mid \mathbf{s}_{t-1}) \, g(\mathbf{v}_t \mid \mathbf{s}_t)}{p(\mathbf{v}_{1:t})} \nonumber \\
&= \frac{p(\mathbf{s}_{1:t-1}, \mathbf{v}_{1:t-1})}{p(\mathbf{v}_{1:t-1})} \frac{f(\mathbf{s}_t \mid \mathbf{s}_{t-1}) \, g(\mathbf{v}_t \mid \mathbf{s}_t)}{p(\mathbf{v}_t \mid \mathbf{v}_{1:t-1})} \nonumber \\
&= p(\mathbf{s}_{1:t-1} \mid \mathbf{v}_{1:t-1}) \frac{f(\mathbf{s}_t \mid \mathbf{s}_{t-1}) \, g(\mathbf{v}_t \mid \mathbf{s}_t)}{p(\mathbf{v}_t \mid \mathbf{v}_{1:t-1})}
\label{eq:recursive_tfd}
\end{align}

We can get the filtering distribution by marginalizing:

\begin{align}
p(\mathbf{s}_t \mid \mathbf{v}_{1:t}) &= \int p(\mathbf{s}_{1:t} \mid \mathbf{v}_{1:t}) \, d\mathbf{s}_{1:t-1} \nonumber \\
&= \int p(\mathbf{s}_{1:t-1} \mid \mathbf{v}_{1:t-1}) \frac{f(\mathbf{s}_t \mid \mathbf{s}_{t-1}) \, g(\mathbf{v}_t \mid \mathbf{s}_t)}{p(\mathbf{v}_t \mid \mathbf{v}_{1:t-1})} \, d\mathbf{s}_{1:t-1} \nonumber \\
&= \frac{g(\mathbf{v}_t \mid \mathbf{s}_t)}{p(\mathbf{v}_t \mid \mathbf{v}_{1:t-1})} \int p(\mathbf{s}_{1:t-1} \mid \mathbf{v}_{1:t-1}) \, f(\mathbf{s}_t \mid \mathbf{s}_{t-1}) \, d\mathbf{s}_{1:t-1} \nonumber \\
&= \frac{g(\mathbf{v}_t \mid \mathbf{s}_t)}{p(\mathbf{v}_t \mid \mathbf{v}_{1:t-1})} \overbrace{\int p(\mathbf{s}_{1:t} \mid \mathbf{v}_{1:t-1}) \, d\mathbf{s}_{1:t-1}}^{= p(\mathbf{s}_t \mid \mathbf{v}_{1:t-1}) \text{ by marginalization}} \nonumber \\
&= \frac{p(\mathbf{s}_t \mid \mathbf{v}_{1:t-1}) \, g(\mathbf{v}_t \mid \mathbf{s}_t)}{p(\mathbf{v}_t \mid \mathbf{v}_{1:t-1})}
\label{eq:sfd_deriv}
\end{align}

This is the same result from the prediction-correction steps. The two most important equations for particle filtering are:

\begin{equation}
\label{eq:TFD}
\boxed{p(\mathbf{s}_{1:t} \mid \mathbf{v}_{1:t}) = p(\mathbf{s}_{1:t-1} \mid \mathbf{v}_{1:t-1}) \frac{f(\mathbf{s}_t \mid \mathbf{s}_{t-1}) \, g(\mathbf{v}_t \mid \mathbf{s}_t)}{p(\mathbf{v}_t \mid \mathbf{v}_{1:t-1})}}
\end{equation}

which we call the ``Trajectory Filtering Distribution'' (TFD), and:

\begin{equation}
\label{eq:SFD}
\boxed{p(\mathbf{s}_t \mid \mathbf{v}_{1:t}) = \frac{p(\mathbf{s}_t \mid \mathbf{v}_{1:t-1}) \, g(\mathbf{v}_t \mid \mathbf{s}_t)}{p(\mathbf{v}_t \mid \mathbf{v}_{1:t-1})}}
\end{equation}

which we call the ``State Filtering Distribution'' (SFD).

\section{Particle Filtering}

\subsection{Basics of Monte Carlo and Importance Sampling}

Most systems aren't linear and/or Gaussian, so we need approximate inference. The Monte Carlo method approximates integrals and expectations with random samples:

\begin{equation}
\label{eq:mc}
\mathbb{E}_{p(\mathbf{x})}[f(\mathbf{x})] = \int f(\mathbf{x}) \, p(\mathbf{x}) \, d\mathbf{x} \approx \frac{1}{N} \sum_{n=1}^{N} f(\mathbf{x}_n) \qquad \mathbf{x}_n \sim p(\mathbf{x})
\end{equation}

These samples can also approximate the target distribution itself:

\begin{equation}
p(\mathbf{x}) \approx \frac{1}{N} \sum_{n=1}^{N} \delta_{\mathbf{x}}(\mathbf{x}_n)
\end{equation}

where $\delta_{\mathbf{x}}(\mathbf{x}_n)$ is the Dirac delta mass evaluated at point $\mathbf{x}_n$.

Often it is not possible to sample from the distribution of interest. We use \textit{importance sampling}, which samples from another known distribution and assigns weights:

\begin{align}
\mathbb{E}_{p(\mathbf{x})}[f(\mathbf{x})] &= \int f(\mathbf{x}) \cdot p(\mathbf{x}) \, d\mathbf{x} \nonumber \\
&= \int \frac{f(\mathbf{x}) \cdot p(\mathbf{x})}{q(\mathbf{x})} \cdot q(\mathbf{x}) \, d\mathbf{x} \nonumber \\
&= \mathbb{E}_{q(\mathbf{x})} \left[ f(\mathbf{x}) \cdot \frac{p(\mathbf{x})}{q(\mathbf{x})} \right] \nonumber \\
&= \mathbb{E}_{q(\mathbf{x})} \left[ f(\mathbf{x}) \cdot w(\mathbf{x}) \right]
\label{eq:importance_sampling}
\end{align}

The distribution $q(\mathbf{x})$ is called the \textcolor{orange}{\textbf{proposal}}, and the weight $w(\mathbf{x}) = p(\mathbf{x})/q(\mathbf{x})$ adjusts for sampling from the wrong distribution.

In Bayesian inference, we have a target posterior $\pi(\mathbf{x}) = p(\mathbf{x} \mid \mathcal{D})$. Since we usually can only evaluate $\pi(\mathbf{x})$ up to a normalizing constant, we use the \textit{self-normalized} IS estimator:

\begin{align}
\mathbb{E}_{\pi}[f(\mathbf{x})] &= \int f(\mathbf{x}) \, \pi(\mathbf{x}) \, d\mathbf{x} \nonumber \\
&= \int f(\mathbf{x}) \frac{\pi(\mathbf{x})}{q(\mathbf{x})} q(\mathbf{x}) \, d\mathbf{x} \nonumber \\
&= \int f(\mathbf{x}) \frac{p(\mathbf{x}, \mathcal{D})}{p(\mathcal{D})q(\mathbf{x})} q(\mathbf{x}) \, d\mathbf{x} \nonumber \\
&= \frac{1}{p(\mathcal{D})} \int f(\mathbf{x}) \frac{p(\mathbf{x}, \mathcal{D})}{q(\mathbf{x})} q(\mathbf{x}) \, d\mathbf{x} \nonumber \\
&= \frac{1}{\mathbb{E}_q\left[\frac{p(\mathbf{x}, \mathcal{D})}{q(\mathbf{x})}\right]} \cdot \mathbb{E}_q\left[f(\mathbf{x}) \frac{p(\mathbf{x}, \mathcal{D})}{q(\mathbf{x})}\right] \nonumber \\
&\approx \frac{1}{\sum_{n=1}^{N} \frac{p(\mathbf{x}_n, \mathcal{D})}{q(\mathbf{x}_n)}} \cdot \sum_{n=1}^{N} f(\mathbf{x}_n) \frac{p(\mathbf{x}_n, \mathcal{D})}{q(\mathbf{x}_n)} := \widehat{\mathcal{I}}_{SN}
\label{eq:self_normalized_is}
\end{align}

The ratio $\tilde{w}(\mathbf{x}_n) := \frac{p(\mathbf{x}_n, \mathcal{D})}{q(\mathbf{x}_n)}$ is the unnormalized importance weight. The approximate posterior is:

\begin{equation}
\pi(\mathbf{x}) \approx \sum_{n=1}^{N} w(\mathbf{x}_n) \delta_{\mathbf{x}_n}(\mathbf{x}) \qquad w(\mathbf{x}_n) = \frac{\tilde{w}(\mathbf{x}_n)}{\sum_{k=1}^{N} \tilde{w}(\mathbf{x}_k)}
\end{equation}

\subsection{Choice of Proposal and Variance of Importance Weights}

The proposal should minimize the variance of our estimators. For the non-normalized estimator:

\begin{equation}
\mathbb{V}_q[\widehat{\mathcal{I}}_{NN}] = \mathbb{E}_q\left[\left(\widehat{\mathcal{I}}_{NN} - \mathcal{I}\right)^2\right]
\end{equation}

Using the variance identity:

\begin{equation}
\mathbb{V}_q[\widehat{\mathcal{I}}_{NN}] = \frac{1}{N} \mathbb{E}_q\left[\left(\frac{f(\mathbf{x})\pi(\mathbf{x})}{q(\mathbf{x})}\right)^2\right] - \frac{1}{N}(\mathcal{I})^2
\end{equation}

Expanding the first term:

\begin{align}
\mathbb{E}_q\left[\left(\frac{f(\mathbf{x})\pi(\mathbf{x})}{q(\mathbf{x})}\right)^2\right] &= \int \frac{(f(\mathbf{x})\pi(\mathbf{x}))^2}{q(\mathbf{x})} \, d\mathbf{x} \nonumber \\
&= \int |f(\mathbf{x})\pi(\mathbf{x})| \frac{|f(\mathbf{x})\pi(\mathbf{x})|}{q(\mathbf{x})} \, d\mathbf{x}
\label{eq:variance_expansion}
\end{align}

The proposal that minimizes variance is:

\begin{equation}
q^{\star}(\mathbf{x}) = \frac{|f(\mathbf{x})\pi(\mathbf{x})|}{\int |f(\mathbf{x})\pi(\mathbf{x})| \, d\mathbf{x}}
\end{equation}

\subsection{Sequential Importance Sampling}

We now sequentially estimate distributions of the form $(p(\mathbf{s}_{1:t} \mid \mathbf{v}_{1:t}))_t$ using weighted samples (particles).

Let $\gamma_t(\mathbf{s}_{1:t})$ be the unnormalized target distribution at time $t$. Let $\pi(\mathbf{s}_{1:t})$ be the normalized version: $\pi(\mathbf{s}_{1:t}) = \gamma_t(\mathbf{s}_{1:t})/Z_t$ with $Z_t = \int \gamma_t(\mathbf{s}_{1:t}) \, d\mathbf{s}_{1:t}$.

Using IS directly with a proposal:

\begin{equation}
\label{eq:sis_weights}
\tilde{w}_t = \frac{\gamma_t(\mathbf{s}_{1:t})}{q_t(\mathbf{s}_{1:t})}
\end{equation}

We can approximate the posterior:

\begin{equation}
p(\mathbf{s}_{1:t} \mid \mathbf{v}_{1:t}) \approx \sum_{n=1}^{N} w_t^n \delta_{\mathbf{s}_{1:t}}(\mathbf{s}_{1:t}^n) \qquad \mathbf{s}_{1:t}^n \sim q_t(\mathbf{s}_{1:t})
\end{equation}

To avoid scaling linearly with dimension $t$, we impose an autoregressive structure on the proposal:

\begin{equation}
q_t(\mathbf{s}_{1:t}) = q_{t-1}(\mathbf{s}_{1:t-1}) \, q_t(\mathbf{s}_t \mid \mathbf{s}_{1:t-1})
\end{equation}

This allows us to express weights at time $t$ recursively:

\begin{align}
\tilde{w}_t(\mathbf{s}_{1:t}) &= \frac{\gamma_t(\mathbf{s}_{1:t})}{q_t(\mathbf{s}_{1:t})} \nonumber \\
&= \frac{1}{q_{t-1}(\mathbf{s}_{1:t-1})} \frac{\gamma_{t-1}(\mathbf{s}_{1:t-1})}{\gamma_{t-1}(\mathbf{s}_{1:t-1})} \frac{\gamma_t(\mathbf{s}_{1:t})}{q_t(\mathbf{s}_t \mid \mathbf{s}_{1:t-1})} \nonumber \\
&= \frac{\gamma_{t-1}(\mathbf{s}_{1:t-1})}{q_{t-1}(\mathbf{s}_{1:t-1})} \frac{\gamma_t(\mathbf{s}_{1:t})}{\gamma_{t-1}(\mathbf{s}_{1:t-1}) \, q_t(\mathbf{s}_t \mid \mathbf{s}_{1:t-1})} \nonumber \\
&= \tilde{w}_{t-1}(\mathbf{s}_{1:t-1}) \cdot \underbrace{\frac{\gamma_t(\mathbf{s}_{1:t})}{\gamma_{t-1}(\mathbf{s}_{1:t-1}) \, q_t(\mathbf{s}_t \mid \mathbf{s}_{1:t-1})}}_{\varpi_t(\mathbf{s}_{t-1}, \mathbf{s}_t)}
\label{eq:sis_weight_update}
\end{align}

where we define the \textit{incremental importance weight} $\varpi_t(\mathbf{s}_{t-1}, \mathbf{s}_t)$.

For our state space model with target $\gamma_t(\mathbf{s}_{1:t}) := p(\mathbf{s}_{1:t}, \mathbf{v}_{1:t})$:

\begin{align}
\varpi_t(\mathbf{s}_{t-1}, \mathbf{s}_t) &= \frac{\gamma_t(\mathbf{s}_{1:t})}{\gamma_{t-1}(\mathbf{s}_{1:t-1}) \, q_t(\mathbf{s}_t \mid \mathbf{s}_{1:t-1}, \mathbf{v}_{1:t})} \nonumber \\
&= \frac{f(\mathbf{s}_t \mid \mathbf{s}_{t-1}) \, g(\mathbf{v}_t \mid \mathbf{s}_t) \, p(\mathbf{s}_{1:t-1}, \mathbf{v}_{1:t-1})}{\gamma_{t-1}(\mathbf{s}_{1:t-1}) \, q_t(\mathbf{s}_t \mid \mathbf{s}_{1:t-1}, \mathbf{v}_{1:t})} \nonumber \\
&= \frac{f(\mathbf{s}_t \mid \mathbf{s}_{t-1}) \, g(\mathbf{v}_t \mid \mathbf{s}_t)}{q_t(\mathbf{s}_t \mid \mathbf{s}_{1:t-1}, \mathbf{v}_{1:t})}
\label{eq:sis_ssm_weights}
\end{align}

SIS has exponentially increasing variance, leading to \textit{weight degeneracy} where one weight $\approx 1$ and all others $\approx 0$.

\subsection{Resampling}

Sequential Monte Carlo (SMC) or Sequential Importance Resampling (SIR) mitigates weight degeneracy by resampling particles. After computing weights, we resample with replacement where each particle is selected with probability equal to its weight, then reset weights to $1/N$.

Resampling can be interpreted as a clever choice of proposal. Instead of propagating particles from $q_{t}(\mathbf{s}_t \mid \mathbf{s}_1)$, we use the approximate posterior from the previous iteration and sample from $\widehat{\pi}_t \cdot q_1(\mathbf{s}_2 \mid \mathbf{s}_1)$.

\subsubsection*{Algorithm 1: Sequential Monte Carlo / Sequential Importance Resampling}

\textbf{At time $t=1$:}
\begin{enumerate}
\item \textbf{Propagation}: sample from proposal $\mathbf{s}_1^n \sim q_1(\mathbf{s}_1)$
\item \textbf{Update}: compute weights as $w_1^n \propto \varpi_1^n(\mathbf{s}_0^n, \mathbf{s}_1^n)$
\item \textbf{Resample}: $(\mathbf{s}_1^n, w_1^n)_{n=1}^N$ to obtain $(\mathbf{r}_1^n, 1/N)_{n=1}^N$
\end{enumerate}

\textbf{At time $t \geq 2$:}
\begin{enumerate}
\item \textbf{Propagation}: sample $\mathbf{s}_t^n \sim q_t(\mathbf{s}_t \mid \mathbf{r}_{1:t-1}^n)$ and set $\mathbf{s}_{1:t}^n \leftarrow (\mathbf{r}_{1:t-1}^n, \mathbf{s}_t^n)$
\item \textbf{Update}: compute weights as $w_t^n \propto \varpi_t^n(\mathbf{s}_{t-1}^n, \mathbf{s}_t^n)$
\item \textbf{Resample}: $(\mathbf{s}_{1:t}^n, w_t^n)_{n=1}^N$ to obtain $(\mathbf{r}_{1:t}^n, 1/N)_{n=1}^N$
\end{enumerate}

If we choose the proposal equal to the transition density, $q_t(\mathbf{s}_t \mid \mathbf{s}_{1:t-1}, \mathbf{v}_{1:t}) = f(\mathbf{s}_t \mid \mathbf{s}_{t-1})$, the weight update simplifies to:

\begin{align}
\varpi_t(\mathbf{s}_{t-1}, \mathbf{s}_t) &= \frac{f(\mathbf{s}_t \mid \mathbf{s}_{t-1}) \, g(\mathbf{v}_t \mid \mathbf{s}_t)}{q_t(\mathbf{s}_t \mid \mathbf{s}_{1:t-1}, \mathbf{v}_{1:t})} \nonumber \\
&= \frac{f(\mathbf{s}_t \mid \mathbf{s}_{t-1}) \, g(\mathbf{v}_t \mid \mathbf{s}_t)}{f(\mathbf{s}_t \mid \mathbf{s}_{t-1})} \nonumber \\
&= g(\mathbf{v}_t \mid \mathbf{s}_t)
\label{eq:bpf_weights}
\end{align}

This gives the \textit{Bootstrap Particle Filter} (BPF).

Resampling introduces \textit{path degeneracy}, where repeated resampling causes particle diversity to collapse, with most particles sharing a single ancestor.

\section{Propagating Particles by Incorporating the Current Measurement}

\subsection{The Effect of Using the Locally Optimal Proposal}

The ``optimal'' or ``locally optimal'' proposal in filtering is:

\begin{equation}
q_t(\mathbf{s}_t \mid \mathbf{s}_{1:t-1}, \mathbf{v}_{1:t}) = p(\mathbf{s}_t \mid \mathbf{s}_{t-1}, \mathbf{v}_t) = \frac{g(\mathbf{v}_t \mid \mathbf{s}_t) \, f(\mathbf{s}_t \mid \mathbf{s}_{t-1})}{\int g(\mathbf{v}_t \mid \mathbf{s}_t) \, f(\mathbf{s}_t \mid \mathbf{s}_{t-1}) \, d\mathbf{s}_t}
\end{equation}

With this proposal, the weight update becomes:

\begin{align}
\varpi_t(\mathbf{s}_{t-1}, \mathbf{s}_t) &= \frac{f(\mathbf{s}_t \mid \mathbf{s}_{t-1}) \, g(\mathbf{v}_t \mid \mathbf{s}_t)}{\frac{f(\mathbf{s}_t \mid \mathbf{s}_{t-1}) \, g(\mathbf{v}_t \mid \mathbf{s}_t)}{p(\mathbf{v}_t \mid \mathbf{s}_{t-1})}} \nonumber \\
&= p(\mathbf{v}_t \mid \mathbf{s}_{t-1})
\label{eq:optimal_weight}
\end{align}

This expression does not depend on $\mathbf{s}_t$, so the conditional variance of the weights is zero. However, this proposal is intractable: sampling from it is as hard as the filtering problem itself, and it requires evaluating the predictive likelihood $p(\mathbf{v}_t \mid \mathbf{s}_{t-1})$.

\subsection{The Auxiliary Particle Filter}

\subsubsection{A First Interpretation: A Standard SMC Algorithm with a Different $\gamma$}

The APF can be interpreted as a standard SMC algorithm where the target is $\gamma_t(\mathbf{s}_{1:t}) = p(\mathbf{s}_{1:t}, \mathbf{v}_{1:t+1})$ instead of $p(\mathbf{s}_{1:t}, \mathbf{v}_{1:t})$.

This target decomposes as:

\begin{align}
\gamma_t(\mathbf{s}_{1:t}) &:= p(\mathbf{s}_{1:t}, \mathbf{v}_{1:t+1}) \nonumber \\
&= \int p(\mathbf{s}_{1:t+1}, \mathbf{v}_{1:t+1}) \, d\mathbf{s}_{t+1} \nonumber \\
&= \int p(\mathbf{s}_{1:t}, \mathbf{v}_{1:t}) \cdot f(\mathbf{s}_{t+1} \mid \mathbf{s}_t) \cdot g(\mathbf{v}_{t+1} \mid \mathbf{s}_{t+1}) \, d\mathbf{s}_{t+1} \nonumber \\
&= p(\mathbf{s}_{1:t}, \mathbf{v}_{1:t}) \int f(\mathbf{s}_{t+1} \mid \mathbf{s}_t) \cdot g(\mathbf{v}_{t+1} \mid \mathbf{s}_{t+1}) \, d\mathbf{s}_{t+1} \nonumber \\
&= p(\mathbf{s}_{1:t}, \mathbf{v}_{1:t}) \cdot \underbrace{p(\mathbf{v}_{t+1} \mid \mathbf{s}_t)}_{\text{predictive likelihood}}
\label{eq:apf_target}
\end{align}

The weight update is:

\begin{align}
\varpi_t(\mathbf{s}_{t-1}, \mathbf{s}_t) &= \frac{\gamma_t(\mathbf{s}_{1:t})}{\gamma_{t-1}(\mathbf{s}_{1:t-1}) \, q_t(\mathbf{s}_t \mid \mathbf{s}_{1:t-1})} \nonumber \\
&= \frac{p(\mathbf{s}_{1:t}, \mathbf{v}_{1:t}) \, p(\mathbf{v}_{t+1} \mid \mathbf{s}_t)}{p(\mathbf{s}_{1:t-1}, \mathbf{v}_{1:t-1}) \, p(\mathbf{v}_t \mid \mathbf{s}_{t-1}) \, q_t(\mathbf{s}_t \mid \mathbf{s}_{1:t-1})} \nonumber \\
&= \frac{f(\mathbf{s}_t \mid \mathbf{s}_{t-1}) \, g(\mathbf{v}_t \mid \mathbf{s}_t) \, p(\mathbf{v}_{t+1} \mid \mathbf{s}_t)}{p(\mathbf{v}_t \mid \mathbf{s}_{t-1}) \, q_t(\mathbf{s}_t \mid \mathbf{s}_{1:t-1})}
\label{eq:apf_weight_update}
\end{align}

To approximate $p(\mathbf{s}_{1:t} \mid \mathbf{v}_{1:t})$ at iteration $t$, we use:

\begin{equation}
\varpi_t(\mathbf{s}_{t-1}, \mathbf{s}_t) = \frac{f(\mathbf{s}_t \mid \mathbf{s}_{t-1}) \, g(\mathbf{v}_t \mid \mathbf{s}_t)}{p(\mathbf{v}_t \mid \mathbf{s}_{t-1}) \, q_t(\mathbf{s}_t \mid \mathbf{v}_t, \mathbf{s}_{t-1})}
\end{equation}

In practice, the predictive likelihood is approximated: $\widehat{p}(\mathbf{v}_t \mid \mathbf{s}_{t-1}) = g(\mathbf{v}_t \mid \boldsymbol{\mu}(\mathbf{s}_t))$ where $\boldsymbol{\mu}(\mathbf{s}_t)$ is some likely value.

A popular choice is $\widehat{p}(\mathbf{v}_t \mid \mathbf{s}_{t-1}) = g(\mathbf{v}_t \mid \boldsymbol{\mu}_t)$ where $\boldsymbol{\mu}_t$ is the mean of $f(\mathbf{s}_t \mid \mathbf{s}_{t-1})$, and $q_t(\mathbf{s}_t \mid \mathbf{v}_t, \mathbf{s}_{t-1}) = f(\mathbf{s}_t \mid \mathbf{s}_{t-1})$, giving:

\begin{equation}
\varpi_t(\mathbf{s}_{t-1}, \mathbf{s}_t) = \frac{g(\mathbf{v}_t \mid \mathbf{s}_t)}{g(\mathbf{v}_t \mid \boldsymbol{\mu}_t)}
\end{equation}

\subsubsection{The Original Interpretation of APF and Marginal Particle Filters}

In Marginal PFs, we perform importance sampling with target $p(\mathbf{s}_t \mid \mathbf{v}_{1:t})$. The target is proportional to:

\begin{align}
p(\mathbf{s}_t \mid \mathbf{v}_{1:t}) &\propto g(\mathbf{v}_t \mid \mathbf{s}_t) \, p(\mathbf{s}_t \mid \mathbf{v}_{1:t-1}) \nonumber \\
&= g(\mathbf{v}_t \mid \mathbf{s}_t) \int f(\mathbf{s}_t \mid \mathbf{s}_{t-1}) \, p(\mathbf{s}_{t-1} \mid \mathbf{v}_{1:t-1}) \, d\mathbf{s}_{t-1} \nonumber \\
&\approx g(\mathbf{v}_t \mid \mathbf{s}_t) \sum_{n=1}^{N} w_{t-1}^n \, f(\mathbf{s}_t \mid \mathbf{s}_{t-1}^n)
\end{align}

This can be rewritten as:

\begin{equation}
g(\mathbf{v}_t \mid \mathbf{s}_t) \sum_{n=1}^{N} w_{t-1}^n \, f(\mathbf{s}_t \mid \mathbf{s}_{t-1}^n) = \sum_{n=1}^{N} w_{t-1}^n \, p(\mathbf{v}_t \mid \mathbf{s}_{t-1}^n) \, p(\mathbf{s}_t \mid \mathbf{s}_{t-1}^n, \mathbf{v}_t)
\end{equation}

We choose a mixture proposal:

\begin{equation}
q_t(\mathbf{s}_t \mid \mathbf{v}_{1:t}, \mathbf{s}_{t-1}) = \sum_{n=1}^{N} w_{t-1}^n \, q_t(\mathbf{s}_t \mid \mathbf{v}_t, \mathbf{s}_{t-1}^n)
\end{equation}

The unnormalized importance weight is:

\begin{equation}
\widetilde{w}_t = \frac{g(\mathbf{v}_t \mid \mathbf{s}_t) \sum_{n=1}^{N} w_{t-1}^n \, f(\mathbf{s}_t \mid \mathbf{s}_{t-1}^n)}{\sum_{n=1}^{N} w_{t-1}^n \, q_t(\mathbf{s}_t \mid \mathbf{v}_t, \mathbf{s}_{t-1}^n)}
\end{equation}

In the original APF formulation, Pitt and Shephard consider the joint:

\begin{equation}
p(n, \mathbf{s}_t \mid \mathbf{v}_{1:t}) \propto g(\mathbf{v}_t \mid \mathbf{s}_t) \, w_{t-1}^n \, f(\mathbf{s}_t \mid \mathbf{s}_{t-1}^n) = w_{t-1}^n \, p(\mathbf{v}_t \mid \mathbf{s}_{t-1}^n) \, p(\mathbf{s}_t \mid \mathbf{s}_{t-1}^n, \mathbf{v}_t)
\end{equation}

where $n$ is an auxiliary variable indexing the mixture. The marginal of the index is:

\begin{equation}
p(n \mid \mathbf{v}_{1:t}) \propto w_{t-1}^n \, p(\mathbf{v}_t \mid \mathbf{s}_{t-1}^n)
\end{equation}

Approximating the predictive likelihood as $g(\mathbf{v}_t \mid \boldsymbol{\mu}_t^n)$, we define the ``simulation weight'':

\begin{equation}
\lambda_t^n \propto w_{t-1}^n \, g(\mathbf{v}_t \mid \boldsymbol{\mu}_t^n)
\end{equation}

The proposal is:

\begin{equation}
q_t(n, \mathbf{s}_t \mid \mathbf{v}_{1:t}) = \lambda_t^n \, q_t(\mathbf{s}_t \mid \mathbf{s}_{t-1}^n, \mathbf{v}_t)
\end{equation}

and the importance weight is:

\begin{equation}
w_t(\mathbf{s}_t, n) \propto \frac{g(\mathbf{v}_t \mid \mathbf{s}_t) \, w_{t-1}^n \, f(\mathbf{s}_t \mid \mathbf{s}_{t-1}^n)}{\lambda_t^n \, q_t(\mathbf{s}_t \mid \mathbf{s}_{t-1}^n, \mathbf{v}_t)}
\end{equation}

\section{The Multiple Importance Sampling Interpretation of PF}

\subsection{The Improved Auxiliary Particle Filter}

[This section would continue with the details of the IAPF and MIS interpretation, following the same pattern]

\section*{References}

\begin{enumerate}
\item Elvira et al. - Improved Auxiliary Particle Filter
\item Doucet et al. - Sequential Monte Carlo Methods
\item Elvira et al. - Numerical results
\item Naesseth et al. - Particle filtering survey
\item Marginal Particle Filters reference
\item Pitt and Shephard - APF original paper
\end{enumerate}

\end{document}
